\documentclass[17pt,colspace=23mm,subcolspace=0mm]{tikzposter}
\geometry{paperwidth=40in,paperheight=30in}
\makeatletter
\setlength{\TP@visibletextwidth}{\textwidth-2\TP@innermargin}
\setlength{\TP@visibletextheight}{\textheight-2\TP@innermargin}
\makeatother

\usepackage{amsmath}
\usepackage{amssymb}
\usepackage[
    backend=biber,
    doi=false,
    eprint=false,
    isbn=false,
    sorting=none,
    style=numeric-comp,
    url=false,
]{biblatex}
\usepackage{bm}
\usepackage{comment}
\usepackage{graphicx}
\usepackage[hidelinks]{hyperref}
\usepackage{lipsum}
\usepackage{multicol}
\usepackage{subcaption}
\usepackage{tikz}
\usepackage{titlesec}

% Typography Tweaks
%
\usepackage{fontspec}

% To use a local font (e.g., https://github.com/alerque/libertinus):
%
\setmainfont[
  Path           = ../../fonts/libertinus/,
  UprightFont    = *-Regular,
  BoldFont       = *-Bold,
  ItalicFont     = *-Italic,
  BoldItalicFont = *-BoldItalic,
  Extension      = .otf,
]{LibertinusSerif}

% Alternatively, if you prefer to use a pre-installed system font:
%
% \setmainfont{Libertinus Serif}

% [Optional] Set root graphics directory for images
%
\graphicspath{{../images/}}

% [Optional] References
%
\addbibresource{refs.bib}

% tikz graphics stuff
%
\tikzposterlatexaffectionproofoff  % remove watermark
\usetikzlibrary{decorations.pathreplacing, arrows.meta, positioning, calc, shapes.geometric}

% Shrink section title padding
%
\titlespacing*{\section} {0pt}{0.25cm}{0.25cm}

% See https://identity.stanford.edu/design-elements/color/primary-colors/
%
\definecolor{customaccent}{RGB}{140,21,21}      % Cardinal red
\definecolor{customprimary}{RGB}{46,45,41}      % Black
\definecolor{customprimarymuted}{RGB}{83,86,90} % Cool gray
\definecolor{customsecondary}{RGB}{255,255,255} % White
\definecolor{customaccentbg}{HTML}{EAEAEA}      % 10% black

% Set and re-define tikzposter theme colors/component styles (see https://texdoc.org/serve/tikzposter/0)
% Try changing the `custom*` values around to see how they affect different parts of the poster.
%
\colorlet{framecolor}{customaccent}
\colorlet{backgroundcolor}{customsecondary}
\colorlet{titlefgcolor}{customsecondary}
\colorlet{titlebgcolor}{customprimary}  % try setting this to customprimary to see a fun gradient effect in the title.
\colorlet{blocktitlefgcolor}{customsecondary}
\colorlet{blocktitlebgcolor}{customaccent}
\colorlet{blockbodyfgcolor}{customprimary}
\colorlet{blockbodybgcolor}{customaccentbg}

% Mix and match the below style options to fit the vibe you want for your poster components.
%
\usetitlestyle{Envelope}  % Default,Basic,Envelope,Wave,VerticalShading,Filled,Empty
\usebackgroundstyle{BottomVerticalGradation}  % Default,Rays,VerticalGradation,BottomVerticalGradation,Empty
\useblockstyle{Basic}  % Default,Basic,Minimal,Envelope,Corner,Slide,TornOut
\useinnerblockstyle{Default}  % Default,Table
\usenotestyle{Sticky}  % Default,Corner,VerticalShading,Sticky

% Title matter
%
\title{The Best Poster Ever: A Template}
\author{Student McStudentFace$^{\dagger *}$, Advisor McAdvisorFace$^{\dagger\ddagger}$}
\date{\today}
\institute{
    ${^\dagger}$Institute for Students, 
    $^\ddagger$Department of Really Great Professors
    \textbar\ The University
}

\begin{document}
    % See https://texdoc.org/serve/tikzposter/0 for more information about `titletoblockverticalspace`.
    %
    \maketitle[titletoblockverticalspace=2.0cm]

    % See https://texdoc.org/serve/tikzposter/0 for more information about how columns and sub-columns are laid out.
    %
    \begin{columns}
        \column{0.7}
        \begin{subcolumns}
            \subcolumn{0.6}
            \block{Background}
            {
                \setcounter{section}{0}
                \section{The "Why"} 
               
                % We use the `multicols` package here to separate text and graphic into a side-by-side layout. The 
                % `subcolumns` blocks don't seem to nest nicely, else we'd use that.
                %
                \begin{multicols}{2}
                
                    \subsection{Why Hasn't This Been Done?} 
                    
                    \begin{itemize}
                        \renewcommand{\labelitemi}{$\rightarrow$}
                        \item{
                            \lipsum[1][1]  % placeholder text
                        }
                        \item{
                            \lipsum[1][2-3]  % placeholder text
                        }
                        \item{
                            \lipsum[1][4-5]  % placeholder text
                        }
                    \end{itemize}
                
                    \subsection{Why Can We Do It?} 
           
                    \lipsum[2][1-8]  % placeholder text
    
                    \begin{tikzfigure}[Background Context Graphic]
                        \includegraphics[width=0.175\textwidth]{example-image-a}
                    \end{tikzfigure}
                \end{multicols}
            }

            \subcolumn{0.4}
            \block{Inspiration}
            {
                \setcounter{section}{1}
                \section{The "What"} 

                \lipsum[3][2-3]  % placeholder text
                
                \subsection{Example}
               

                \lipsum[4-5]  % placeholder text
            }
        \end{subcolumns}

        \begin{subcolumns}
            \subcolumn{0.4}
            \block{Theory}
            {
                \setcounter{section}{3}
                \section{The "When"}
    

                \lipsum[6-9]  % placeholder text
            }

            \subcolumn{0.6}
            \block{Implementation}
            {
                \setcounter{section}{4}
                \section{The "How" (Redux)}
    
                \subsection{What Worked?}

                \lipsum[10]  % placeholder text
    
                \subsection{Why Did That Work?}

                \lipsum[11-13]  % placeholder text
            }
            \block{References}
            {
                \nocite{*}
                \printbibliography[heading=none]
            }
        \end{subcolumns}
    
        \column{0.3}
        \begin{subcolumns}
            \subcolumn{1.0}
            \block{Approach}
            {
                \setcounter{section}{2}
                \section{The "How"}
        
                \lipsum[14][1-3]  % placeholder text

                \vspace{0.5cm}
        
                \begin{center}
                    \begin{tikzpicture}[node distance=1.5cm, >=Stealth]
                        % LHS block
                        %
                        \begin{scope}[shift={(0,-1)}]
                            % Side face
                            %
                            \draw[fill=gray!30, draw=black, thick] (3,0) -- (3.5,0.5) -- (3.5,1.5) -- (3,1) -- cycle;
                            
                            % Top face
                            %
                            \draw[fill=gray!10, draw=black, thick] (0,1) -- (3,1) -- (3.5,1.5) -- (0.5,1.5) -- cycle;
                            
                            % Front face
                            %
                            \draw[fill=gray!20, draw=black, thick] (0,0) -- (3,0) -- (3,1) -- (0,1) -- cycle;
                    
                            % Points
                            %
                            % (front)
                            \fill[black] (0.5,0.5) circle (2pt);
                            \fill[black] (1.5,0.3) circle (2pt);
                            \fill[black] (2.5,0.7) circle (2pt);
                            % (top)
                            \fill[black] (1, 1.3) circle (2pt);
                            \fill[black] (2, 1.2) circle (2pt);
                            
                            \node[black, font=\small] at (1.5, -0.5) {From this...};
                        \end{scope}
                    
                        \draw[->, line width=1.5pt, black] (5.5, 0) -- (10.5, 0);
                    
                        % RHS graph
                        %
                        \begin{scope}[shift={(12,0)}]
                            \draw[black, thin] (0,0) -- (1,1) -- (2,0) -- (1,-1) -- cycle;
                            \draw[black, thin] (1,1) -- (1,-1);
                            \draw[black, thin] (0,0) -- (2,0);
                    
                            \foreach \n in {(0,0), (1,1), (2,0), (1,-1), (1,0)} {
                                \fill[black] \n circle (3pt);
                            }
                            
                            \node[black, font=\small] at (1, -1.5) {...to this!};
                        \end{scope}
                    \end{tikzpicture}
                \end{center}
               
                \lipsum[15]  % placeholder text
            }
    
            \block{Results}
            {
                \setcounter{section}{5}
                \section{What We're Publishing}

                \vspace{0.5cm}
    
                \lipsum[16]1-4]  % placeholder text
    
                \begin{tikzfigure}[How we hoped it'd go]
                    \includegraphics[width=0.115\textwidth]{example-image-b}
                \end{tikzfigure}
    
                \lipsum[16][5-9]  % placeholder text
    
                \begin{tikzfigure}[How it really went]
                    \includegraphics[width=0.115\textwidth]{example-image-c}
                \end{tikzfigure}

                \lipsum[16][8-10]  % placeholder text
            }
        \end{subcolumns}

    \end{columns}

    % Footer
    %
    \coordinate (topleft)     at (topright -| bottomleft);
    \coordinate (bottomright) at (topright |- bottomleft);
    \node [
       above right,
       outer sep=0pt,
       draw=none,
       text=customsecondary
    ] at ([shift={(1.6cm,1.0cm)}]bottomleft) {
        \ttfamily\large
        Research Grant McGrant No. 123456
    };

    \node [
       above left,
       outer sep=0pt,
       draw=none,
       text=customsecondary
    ] at ([shift={(-1.6cm,1.0cm)}]bottomright) {
        \large\ttfamily 
        $^*$Corresponding author: 

        student@university.edu
    };
\end{document}